\documentclass[11pt,letterpaper]{article}
\usepackage[utf8]{inputenc}
\usepackage[T1]{fontenc}
\usepackage{lmodern}
\usepackage[margin=0.28in]{geometry}
\usepackage{hyperref}
\hypersetup{colorlinks=false, urlcolor=black, pdfborder=0 0 0}

\setlength{\parindent}{0pt}
\setlength{\parskip}{0pt}
\renewcommand{\baselinestretch}{0.82}
\setlength{\lineskip}{0pt}
\setlength{\lineskiplimit}{0pt}

\makeatletter
\renewcommand{\@listi}{%
  \leftmargin=\leftmargini
  \rightmargin=0pt
  \labelsep=4pt
  \labelwidth=16pt
  \itemindent=0pt
  \listparindent=0pt
  \topsep=0pt
  \partopsep=0pt
  \parsep=0pt
  \itemsep=-11pt
}
\let\saveditemize\itemize
\let\savedenditemize\enditemize
\renewcommand{\itemize}{\bgroup\setlength{\baselineskip}{0.65\baselineskip}\saveditemize}
\renewcommand{\enditemize}{\savedenditemize\egroup}
\makeatother

\begin{document}

{\Large \textbf{HSU, Ke-Cheng (Clio)}}\\[1pt]
(852) 6840 9753 \quad $|$ \quad \href{mailto:khsu@connect.ust.hk}{khsu@connect.ust.hk} \quad $|$ \quad \href{https://github.com/cliohsu1997}{github.com/cliohsu1997} \quad $|$ \quad \href{https://cliohsu1997.github.io/}{cliohsu1997.github.io}

\vspace{0.15pt}

{\large \textbf{Professional Summary}}\\[-0.5pt]
PhD candidate in Economics at HKUST (expected June 2027) seeking Haitong International Young Talent Program 2027 summer internship---preferably \textbf{Research} or cross-border / China-market desks. Strong empirical analysis on large datasets; clear English and Mandarin writing and presentation; works to deadlines with a coachable project style. Hong Kong permanent resident; available for the June 2027 10-week program.

\vspace{0.15pt}

{\large \textbf{Education}}\\[-0.5pt]
\textbf{The Hong Kong University of Science and Technology}, Hong Kong\\
\textit{Doctor of Philosophy in Economics} \hfill Expected Graduation: June 2027\\
Research interests: Empirical industrial organization and digital economics.\\
\textit{Master of Philosophy in Economics} \hfill 2020--2022\\
CGA: 4.1/4.3.\\
\textit{Bachelor of Science in Mathematics and Economics, Minor in Actuarial Mathematics} \hfill 2015--2019\\
CGA: 3.9/4.3.\\
Awards: Dean's List (2016--2019); University Scholarship (730K HKD); Academic Achievement Medal (top 1\%); Postgraduate studentship (1.5M HKD); Redbird scholarship (40K HKD).

\vspace{0.15pt}

{\large \textbf{Experience}}\\[-0.5pt]
\textbf{Research Assistant}, National Taiwan University \hfill 2020\\
\vspace{-6pt}
\begin{itemize}
  \item Modeled crisis intervention as regulator--creditor negotiation; derived and simulated equilibrium predictions to support policy design.
\end{itemize}
\textbf{Teaching Assistant \& Lecturer}, Department of Economics, HKUST \hfill 2020--present\\
\vspace{-6pt}
\begin{itemize}
  \item Lecturer for MSc Math Camp (4.53/5); TA for PhD and undergraduate economics (avg 4.72/5); present technical material clearly in English and Mandarin; Teaching Assistant Coordinator, Business School (2025--2026).
\end{itemize}

\vspace{0.15pt}

{\large \textbf{Projects}}\\[-0.5pt]
\textbf{Inventory and Pricing under Resale Price Maintenance (RPM)} \textit{(w/ K.\ Kawaguchi, HKUST; K.\ Kawai, U.\ Tokyo)}\\
\vspace{-6pt}
\begin{itemize}
  \item Built a dynamic oligopoly model of pricing and inventory under demand uncertainty and RPM, with equilibrium stockouts and unsold returns (35--40\% of Japanese books printed are discarded); constructed title-by-market-by-month panels from monthly bookstore POS data (deliveries, sales, returns; tens of billions of records across 1{,}000+ stores), linked via standardized ISBNs and merged with local demographics and detailed title-level information scraped from Amazon in Python.
  \item Estimated demand and supply with instrumental variables (2SLS) and Hausman-style price instruments; validated the estimator with Monte Carlo on simulated equilibria; selected controls with LASSO and panel fixed effects; implemented iterative equilibrium simulation in C++.
  \item Compared the dynamic model against the classical static model and showed the static version overstates marginal cost by \(\sim\)20\% once demand uncertainty and inventory decisions are ignored; uncovered stockout-driven intertitle and intertemporal substitution (not interregional) through weekly descriptive statistics, regressions, and figures.
\end{itemize}
\vspace{-3pt}
\textbf{Return Policy Design in E-Commerce} \textit{(Sole-authored Job Market Paper)}\\
\vspace{-6pt}
\begin{itemize}
  \item Sole-authored Ph.D.\ structural industrial-organization study: built and estimated a model of purchase, return-shipping-insurance take-up, and return under monopoly pricing, using within-retailer switches between automatic provision and selective adoption of return-shipping insurance on Taobao (2017--2023; \(\sim\)62M RMB revenue; six apparel categories; \(\sim\)20--30\% return rates); scraped and OCR-digitized product attributes in Python.
  \item Developed an original welfare framework that decomposes losses into ex-ante purchase inefficiency and ex-post return inefficiency; recovered category-specific demand, fit uncertainty, and return responses via GMM with bootstrap inference; built a reproducible estimation pipeline on AWS and HPC/SLURM (Python, R, C++) with Monte Carlo validation of policy rankings for seller profit, consumer surplus, and social welfare.
  \item Documented a purchase--return efficiency trade-off and showed that a menu pairing a discounted price with a moderate (\(\sim\)10--15\%) restocking fee raises social welfare by about 2 million RMB (\(\sim\)5\%) and seller profit by about 1.2 million RMB relative to status-quo insurance policies, attaining \(\sim\)75--78\% of first-best welfare outside dresses; automatic provision expands valuable purchases in high-margin categories while selective adoption outperforms it by \(\sim\)15--30 pp of first-best welfare in thinner-margin categories.
\end{itemize}

\vspace{0.15pt}

{\large \textbf{Skills}}\\[-0.5pt]
\textbf{Languages:} Mandarin (Native), English (Fluent), Cantonese (Basic).\\
\textbf{Technical:} Python, R, Stata, C++, Microsoft Excel, Word, PowerPoint, \LaTeX, AWS, HPC/SLURM; MLE, GMM, bootstrap; research reporting and presentations.\\
\textbf{References:} Kei Kawai (UTokyo; Compass Lexecon), Kohei Kawaguchi (HKUST; UTokyo Econ Consulting), Pak Hung Au (HKUST).

\end{document}
